\par\noindent\begin{minipage}{\linewidth}
% !TeX root = ../教材题库_全部条目_转写工作稿.tex
% Source: XJ-B1-C04-PR-P103-Y1012
% 第4章 / 4.1 实数指数幂和幂函数 / 教材 p.103 / PDF 第110页
% 已对照原页转写题干及公式；待教师复核；未制作答案
% 原题及相邻原解截图（核对用）：../题目原图/XJ-B1-C04-PR-P103-Y1012-01.webp
% 本文件只在转写完成后作为题干使用；OCR 条目尚不是可靠题干。
\begin{bankitem}{p.103 练习 第2题}{第4章 / 4.1 实数指数幂和幂函数 / 教材 p.103 / PDF 第110页}{XJ-B1-C04-PR-P103-Y1012}{已对照原页转写题干及公式；待教师复核；未制作答案}
根据幂函数 $y=x^\alpha$ 的图象填空：
\subq{1}当 $\alpha>1$，$x\ge0$ 时，图象恒过点\blank 和\blank；当 $0<x<1$ 时，图象在直线 $y=x$ 的\blank 方；当 $x>1$ 时，图象在直线 $y=x$ 的\blank 方。
\subq{2}当 $0<\alpha<1$，$x\ge0$ 时，图象恒过点\blank 和\blank；当 $0<x<1$ 时，图象在直线 $y=x$ 的\blank 方；当 $x>1$ 时，图象在直线 $y=x$ 的\blank 方。
\subq{3}当 $\alpha<0$，$x>0$ 时，图象恒过点\blank。
\end{bankitem}
\end{minipage}\par\medskip
